\chapter{Classification and Cohomology of $C_2$-Modules}\label{appendix:classandcohomC2modules}

In this appendix we first recall some basic facts about the cohomology of $C_2$-modules. We then give a classification theorem for finitely-generated $C_2$-modules, which reduces the classification of arbitrary $C_2$-modules to that of $C_2$-modules which are $2$-power torsion, a complicated but ultimately algorithmically finite problem.


\section{Cohomology of $C_2$-modules}\label{sect:cohomC2modules}
Write $C_2 = \langle \epsilon \rangle$ for the group with two elements, where $\epsilon^2 = 1$. A $C_2$-module is just a module over the group algebra $\bZ [C_2]$, i.e. an additive group $A=(A,+)$ with an order two automorphism. 
To describe the cohomology of a $C_2$-module $A$, we need invariants and skew-invariants
\[
A^0 \coloneqq \{ x\in A \;\mid\; \epsilon x = x \}, \ \
A^1 \coloneqq \{ x\in A \;\mid\; \epsilon x = -x \}.
\]
\begin{prop} \label{prop:C2cohomology}
  The normalised bar resolution of the trivial $\bZ[C_2]$-module $\bZ$ is the free resolution
  \[
  0 \leftarrow \bZ \xleftarrow{\emph{\text{aug}}} \bZ[C_2] \xleftarrow{\epsilon - 1} \bZ[C_2] \xleftarrow{\epsilon + 1}  \bZ[C_2] \xleftarrow{\epsilon - 1} \bZ[C_2] \xleftarrow{\epsilon + 1} \cdots .
  \]
In particular, the cohomology groups of a $C_2$-module $A$ are as follows:
\[
\HH^{n}(A) \coloneqq \HH^{n}(C_2, A) =
\begin{cases}
  A^0, & \mbox{ if } n=0, \\
  A^0/(\epsilon + 1) A, & \mbox{ if } n \mbox{ is even and } n>0, \\
  A^1/(\epsilon - 1) A, & \mbox{ if } n \mbox{ is odd}.
\end{cases}
\]
\end{prop}

\begin{proof}
  This is \cite[Ex.\ 6.5.2(4)]{Weib} and \cite[Thm.\ 6.2.2]{Weib}.
\end{proof}

\begin{remark}\label{rem:connectingmaps}
  Suppose that $0 \rightarrow A \xrightarrow{\iota} B \rightarrow C \rightarrow 0$ is a short exact sequence of $C_2$-modules. From the explicit resolution of Proposition \ref{prop:C2cohomology} we can also explicitly describe the connecting homomorphisms in the associated long exact sequence on cohomology. We have
\[
\left\{
\begin{array}{r@{}l@{\quad}r@{}l@{\qquad}l}
  \delta_0      & \colon C^0 \rightarrow A^1/(\epsilon-1)A,               & \: \delta_0(c)    & = \iota^{-1}(\epsilon(b)-b), & \text{if } n=0, \\
  \delta_n      & \colon C^0/(\epsilon+1)C \rightarrow A^1/(\epsilon-1)A, & \: \delta_n(c)    & = \iota^{-1}(\epsilon(b)-b), & \text{if } n \text{ is even and } n>0, \\
  \delta_n      & \colon C^1/(\epsilon-1)C \rightarrow A^0/(\epsilon+1)A, & \: \delta_n(c)    & = \iota^{-1}(\epsilon(b)+b), & \text{if } n \text{ is odd},
\end{array} \right.
\]
  where in each case $b \in B$ is any lift of $c \in C$ to $B$.
\end{remark}

Recall that a group is called an elementary 2-group if it is a direct sum of copies of $\bZ/2\bZ$ (and so any such group is naturally viewed as a vector space over the field $\bF_2$ of two elements).
\begin{cor}\label{cor:cohomisvectorspace}
  The group $\HH^n(A)$ is an elementary 2-group for all $n\geq 1$.
\end{cor}
\begin{proof}
If $(1 \pm \epsilon)x=0$, then $(1 \mp \epsilon)x = 2x$, so $2x=0$ in $\HH^n(A)$ for all $x$.     
\end{proof}  

Let us now compute explicitly the cohomology of some classes of $C_2$-modules.
Let $B$ be an abelian group. We can create three $C_2$-modules from $B$:
\[
B_+ = B, \ \epsilon x = x; \qquad B_- = B, \ \epsilon x = -x; \qquad   B_d = B\oplus B, \ \epsilon (x,y) = (y,x). 
\]
For a prime $p$, we also consider the torsion subgroups
\[
B_p = \{ x\in B \,\mid\, px =0\}, \ \
B_{(p)} = \{ x\in B \,\mid\, p^nx =0 \text{ for some $n > 0$}\}.
\]
For these modules, the answer immediately follows from Proposition~\ref{prop:C2cohomology}:
\begin{equation}\label{eqn:cohomofspecialmodules}
\HH^{n}(B_+) =
\begin{cases}
  B, \\
  B/2B, \\
  B_2,
\end{cases}
\HH^{n}(B_-) =
\begin{cases}
  B_2, \\
  B_2, \\
  B/2B, 
\end{cases}
\HH^{n}(B_d) =
\begin{cases}
  B, & \mbox{ if } n=0, \\
  0, & \mbox{ if } n=2k >0, \\
  0, & \mbox{ if } n=2k+1 .
\end{cases}
\end{equation}

An abelian group $B$ admits division by 2 if  $x\mapsto 2x$ is an automorphism of $B$. We denote its inverse by $x\mapsto x/2$. 
\begin{lemma} \label{cohom_div_by2}
  Suppose a $C_2$-module $A$ admits division by $2$. Then $A\cong (A^0)_+ \oplus (A^1)_-$ as $C_2$-modules and
  $\HH^n(A)=0$ for all $n\geq 1$.  
\end{lemma}
\begin{proof} Clearly,  $A^0 + A^1 = A^0 \oplus A^1$ is a $C_2$-submodule of $A$. The inverse to the embedding map is given by
  $x \mapsto ((\epsilon(x) + x)/2, (\epsilon(x) - x)/2 )$. The second statement follows directly from the formulas (\ref{eqn:cohomofspecialmodules}), noting that if $B$ admits division by $2$ then both $B / 2 B = 0$ and $B_2 = 0$.
\end{proof}  

In our classification, we will be interested in $C_2$-modules of the following form. Let $B$ be a $\bZ[C_2]$-module, and consider a pair of abelian group homomorphisms 
  \[
  \Theta = (\Theta_+, \Theta_-), \ \ \Theta_+ : \bZ^n \rightarrow B , \ \  \Theta_- : \bZ^m \rightarrow B .
  \]

\begin{lemma} \label{even_part}
  Consider an operator
  \[\epsilon (x,y,z) = (\epsilon x +\Theta_+ (y) + \Theta_- (z), y, -z)\]
  on the abelian group $B_\Theta (\bZ) \coloneqq B \oplus \bZ^n \oplus \bZ^m$. 
It defines a $\bZ [C_2]$-module structure if and only if 
\newline
$\Theta_+ (\bZ^n ) \subseteq B^1 = \HZ^1 (B)$ and  $\Theta_- (\bZ^m) \subseteq B^0=\HZ^2(B)$.
\end{lemma}
\begin{proof} We need to have $\epsilon^2=1$. This means  
  \[(x,y,z) = \epsilon^2 (x,y,z) = (x + (\epsilon + 1) \Theta_+ (y) + (\epsilon - 1)\Theta_- (z), y, z)\]
  for all $x,y,z$. This proves the lemma.
\end{proof}

Whenever $B_\Theta (\bZ)$ is a $C_2$-module, $B_{\Theta}(\bZ)$ can be naturally viewed as an extension
\[
0 \rightarrow B \rightarrow B_\Theta (\bZ) \rightarrow \bZ^n_+ \oplus \bZ^m_- \rightarrow 0
\]
of $\bZ^n_+ \oplus \bZ^m_-$ by $B$. By definition, two $C_2$-modules $B_\Theta(\bZ)$ and $B_{\Psi}(\bZ)$ define equivalent extensions when there is an isomorphism $\gamma \colon B_\Theta(\bZ) \rightarrow B_{\Psi}(\bZ)$ such that the diagram
\begin{equation} \label{diagram_BTheta}
  \begin{tikzcd}
	0 & B & {{B}_\Theta(\bZ)} & {\bZ^n_+\oplus \bZ^m_-} & 0 \\
	0 & B & {{B}_\Psi(\bZ)} & {\bZ^n_+\oplus \bZ^m_-} & 0
	\arrow[from=1-1, to=1-2]
	\arrow[from=1-2, to=1-3]
	\arrow["\id"', from=1-2, to=2-2]
	\arrow[from=1-3, to=1-4]
	\arrow["\gamma"', from=1-3, to=2-3]
	\arrow[from=1-4, to=1-5]
	\arrow["\id"', from=1-4, to=2-4]
	\arrow[from=2-1, to=2-2]
	\arrow[from=2-2, to=2-3]
	\arrow[from=2-3, to=2-4]
	\arrow[from=2-4, to=2-5]
\end{tikzcd}\end{equation}
commutes.

\begin{cor} \label{even_part_ext_new}
Two $C_2$-modules $B_\Theta(\bZ)$ and $B_{\Psi}(\bZ)$ define equivalent extensions if and only if $\Theta_+-\Psi_+$ is valued in $(\epsilon - 1)(B) = B^1(C_2, B)$ and $\Theta_--\Psi_-$ is valued in $(\epsilon + 1)(B) = B^2(C_2, B)$, and any extension of $
\bZ^n_+ \oplus \bZ^m_-$ by $B$ is of the form $B_\Theta(\bZ)$. In particular,
\[\Ext_{\bZ[C_2]}^1 (\bZ^n_+\oplus \bZ^m_-, B) \cong \Hom (\bZ^n,\HH^1 (B)) \oplus \Hom (\bZ^m,\HH^2(B))\, .\] 
\end{cor}  

\begin{proof}
  From the definition, any equivalence has to be of the form $\gamma(x,y,z) = (x + u_+(y) + u_-(z), y,z)$ for some $u_+ \in \Hom(\bZ^n, B)$ and $u_- \in \Hom(\bZ^m, B)$, and the requirement that $\gamma$ commutes with the $C_2$-action is exactly that $\Theta_+-\Psi_+ = (\epsilon - 1)u_+$ and $\Theta_- -\Psi_- = (\epsilon + 1)u_-$. Conversely, suppose that $(\Theta_+-\Psi_+)$ is valued in $(\epsilon - 1)(B)$ and $(\Theta_--\Psi_-)$ is valued in $(\epsilon + 1)(B)$. Letting $e_1, ... , e_n$ denote the standard basis of $\bZ^n$, one can define $u_+ \in \Hom(\bZ^n, B)$ with $\Theta_+-\Psi_+ = (\epsilon - 1)u_+$ by setting $u_+(e_i) = b_i$ for $i = 1,...,n$, where $(\Theta_+-\Psi_+)(e_i) = (\epsilon - 1)(b_i)$. Similarly, one obtains $u_+ \in \Hom(\bZ^n, B)$ with $\Theta_+-\Psi_+ = (\epsilon - 1)u_+$, and therefore $\gamma(x,y,z) = (x + u_+(y) + u_-(z), y,z)$ defines the required isomorphism.
  
  For the final claim, note that any $C_2$-module structure on $B \oplus \bZ^n \oplus \bZ^m$ for which the exact sequence
  \[
  0 \rightarrow B \rightarrow B \oplus \bZ^n \oplus \bZ^m \rightarrow \bZ^n_+ \oplus \bZ^m_- \rightarrow 0
  \]
  is a sequence of $C_2$-modules must be of the form $B_{\Theta}(\bZ)$ for some $\Theta = (\Theta_+, \Theta_-)$.
\end{proof}

We can describe the cohomology of the modules $B_{\Theta}(\bZ)$, when $B$ is finite, in terms of the maps
\begin{equation}\label{BTheta_reduced}
\overline{\Theta_+} : \bF_2^n = (\bZ/2\bZ)^n  \rightarrow \HH^1(B) , \ \  \overline{\Theta_- } : \bF_2^m = (\bZ/2\bZ)^m \rightarrow \HH^2 (B).
\end{equation}

\begin{prop}\label{prop:cohomofextensionexplicit} Suppose that $B$ is a finite $C_2$-module. Then
\[
\HH^{0}(B_\Theta (\bZ)) \cong \HH^0 (B) \oplus \bZ^n,
\]
and for $r \geq 1$,
\[
 \dim_{\bF_2} \HH^r (B_\Theta (\bZ)) = 
\begin{cases}
\dim_{\bF_2} \HH^r(B) + n - d , & \mbox{ if $r$ is even}, \\
\dim_{\bF_2} \HH^r(B) + m - d, & \mbox{ if $r$ is odd},
\end{cases}
\]
where $d \coloneqq \rank_{\bF_2} \overline{\Theta_+} + \rank_{\bF_2} \overline{\Theta_-}$.
\end{prop}

\begin{proof}
  Considering the long exact sequence on cohomology, one can use Remark \ref{rem:connectingmaps} to explicitly compute the connecting maps as 
  \[
  \begin{tikzcd}
  0 \rar & \HH^0(B) \rar & \HH^0(B_{\Theta}(\bZ)) \rar & \bZ^n \ar[out=-20, in=160, "\Theta_+"]{dll} \\
  & \HH^1(B) \rar & \HH^1(B_{\Theta}(\bZ)) \rar & \bF_2^m \ar[out=-20, in=160, "\overline{\Theta_-}"]{dll} \\
  & \HH^2(B) \rar & \HH^2(B_{\Theta}(\bZ)) \rar & \bF_2^n \ar[out=-20, in=160, "\overline{\Theta_+}"]{dll}\\
  & \HH^3(B) \rar & \cdots
\end{tikzcd}
\]
We have that $\ker(\Theta_+ \colon \bZ^n \rightarrow \HH^1(B)) \cong \bZ^n$, as $\HH^1(B)$ is finite, and $\bZ^n$ is free, hence $\HH^{0}(B_\Theta (\bZ)) \cong \HH^0 (B) \oplus \bZ^n$. For the remaining formula, note that all terms from $\HH^1(B)$ onwards in the sequence are $\bF_2$-vector spaces by Corollary \ref{cor:cohomisvectorspace}. Therefore if $r$ is odd, we have that 
\[
\HH^r(B_{\Theta}(\bZ)) \cong \HH^r(B) / \overline{\Theta_+}(\bF_2^n) \oplus \ker(\overline{\Theta_-} \colon \bF_2^m \rightarrow \HH^{r+1}(B)),
\]
and thus $\dim_{\bF_2} \HH^r (B_\Theta (\bZ)) = \HH^r(B) - \dim_{\bF_2}(\overline{\Theta_+}(\bF_2^n)) + m - \rank_{\bF_2}(\overline{\Theta_-})$. Similarly for even $r$.
\end{proof}

\section{Classification of $C_2$-modules}
In this section give a classification theorem for finitely-generated $C_2$-modules.
As part of our classification, we show that any such module decomposes uniquely into a direct sum of indecomposable modules, i.e., the Krull-Schmidt Theorem holds.
This is mildly surprising since the Krull-Schmidt Theorem fails for finitely generated $\bZ [G]$-modules for a finite group $G$ in general \cite{ReinerKS}.

We show that the indecomposable $C_2$-modules fall into three categories: $p$-torsion for odd $p$, torsion-free and mixed.
The latter are the modules $B_\Theta (\bZ)$, where $B$ is non-trivial and has order a power of $2$, which have 2-power-torsion and torsion-free parts.
Their classification is not explicit, reducing to the classification of Levy structures, which we introduce below.

\subsection{Torsion-free modules}

We first consider $C_2$-modules which are finite-free as $\bZ$-modules. Examples are the modules $\bZ_d$, $\bZ_+$ and $\bZ_-$, from Section \ref{sect:cohomC2modules}. The following shows that these suffice to describe all such modules.

\begin{prop}\label{prop:torsionfreeKS}
  The $C_2$-modules $\bZ_d$, $\bZ_+$ and $\bZ_-$ are indecomposable, and any finitely generated $C_2$-module $M$ which is free as a $\bZ$-module is of the form
  \[
    M \cong \bZ_d^{r_d} \oplus \bZ_{+}^{r_+} \oplus \bZ_{-}^{r_-}
  \]
  for unique $r_d, r_+, r_- \in \bZ_{\geq 0}$. 
\end{prop}

More generally, one can classify finitely generated $C_p$-modules which are free over $\bZ$, where $p$ is any prime \cite[Thm.\ 74.3, Ex.\ 74(4)]{CR0}. We give a proof of Proposition \ref{prop:torsionfreeKS}, which is much simpler than the proof in the general case. We also note that Proposition \ref{prop:torsionfreeKS} is equivalent to a classification of self-inverse integer matrices, which can be done in a very elementary, but again longer, way \cite{Han}.

\begin{proof}
It is direct to see that $\bZ_+$ and $\bZ_-$ indecomposable, this being true even as $\bZ$-modules. For $\bZ_d$, note that $\bZ_d \cong \bZ[C_2]$ as $\bZ[C_2]$-modules, and the reduction modulo $2$ is $\bF_2[C_2]$. Therefore any decomposition $M_1 \oplus M_2 = \bZ_d$ as $\bZ[C_2]$-modules induces a decomposition $\overline{M_1} \oplus \overline{M_2} = \bF_2[C_2]$ of the indecomposable $\bF_2[C_2]$-module $\bF_2[C_2]$, and therefore one of $M_1$ and $M_2$ must have $\bZ$-rank $2$. Both $M_1$ and $M_2$ are free over $\bZ$, being submodules of a free module, hence one of $M_1$ and $M_2$ must be $0$.

Now suppose that $M$ is a arbitrary $C_2$-module, free over $\bZ$, and consider $X \coloneqq M / M^1$. We first show that $X^1 = 0$. Suppose that $\epsilon(x) = - x$, and let $m \in M$ be a lift of $x$, so $\epsilon(m) = - m + m'$ for some $m' \in M^1$. Then $\epsilon(m') = -m'$, but also
\[
  \epsilon(m') = \epsilon(\epsilon(m) + m) = \epsilon(m) + m = m',
\]
so $2m' = 0$, hence $m' = 0$, as $M$ is torsion-free. In particular, $m \in M^1$, so $x = 0$ in $X$. This shows that $X^1 = 0$, from which it follows that $X^0 = X$: for any $x \in X$, $\epsilon(x) = x$, because
\[
\epsilon(\epsilon(x) - x) = x - \epsilon(x) = - (\epsilon(x) - x).
\]
In summary, $M^1 \cong \bZ_{-}^{a}$ and $X \cong \bZ_{+}^{b}$ for some $a, b \geq 0$, and $M$ fits into a short exact sequence
\[
0 \rightarrow \bZ_{-}^{a} \rightarrow M \rightarrow \bZ_{+}^{b} \rightarrow 0.
\]

Choosing a splitting of $M \rightarrow \bZ_{+}^{b}$ as a $\bZ$-module provides a basis $e_1, ... , e_a, f_1, ... , f_b$ of $M$ for which
\[
\begin{pmatrix}
  -I_a & U \\ 0 & I_b
\end{pmatrix}
\]
is the matrix of $\epsilon$. If $P \in \GL_a(\bZ)$ and $Q \in \GL_b(\bZ)$, then
\[
\begin{pmatrix}
  P & 0 \\ 0 & Q^{-1}
\end{pmatrix}
\begin{pmatrix}
  -I_a & U \\ 0 & I_b
\end{pmatrix}
\begin{pmatrix}
  P^{-1} & 0 \\ 0 & Q
\end{pmatrix} = \begin{pmatrix}
  -I_a & PUQ \\ 0 & I_b
\end{pmatrix},
\]
and therefore by changing basis we may assume that $U$ is in Smith normal form. We therefore have that $M$ is the direct sum of the $C_2$-submodules $\langle e_i, f_i \rangle$ for $1 \leq i \leq \min(a,b)$ and either modules $\langle e_i \rangle \cong \bZ_-$ for $i > b$, if $a>b$, or modules $\langle f_i \rangle \cong \bZ_+$ for $i>a$, if $a<b$. It therefore remains to understand the submodules $M_i = \langle e_i, f_i \rangle$ for $1 \leq i \leq \min(a,b)$, where $\epsilon(e_i) = -e_i$ and $\epsilon(f_i) = u_{ii} e_i + f_i$.

If $u_{ii} = 2k$ is even, then with respect to the basis $e_i' \coloneqq e_i, f_i' \coloneqq f_i + k e_i$, $\epsilon(e_i') = - e_i'$ and
\[
\epsilon(f_i') = \epsilon(f_i + k e_i) = u_{ii} e_i + f_i - k e_i = f_i + k e_i = f_i',
\]
hence $M_i \cong \bZ_+ \oplus \bZ_-$.

If instead $u_{ii} = 2k + 1$ is odd, then with respect to the basis $e_i' \coloneqq e_i + (f_i + k e_i), f_i' \coloneqq f_i + k e_i$,
\begin{align*}
\epsilon(f_i') &= \epsilon(f_i + k e_i) = u_{ii} e_i + f_i - k e_i = (f_i + k e_i) + e_i = e_i', \\
\epsilon(e_i') &= \epsilon((f_i + k e_i) + e_i) = e_i' - e_i = f_i + k e_i = f_i',
\end{align*}
and so $M_i \cong \bZ_d$.

To see that this description is unique, suppose that 
\[
\bZ_d^{r_d} \oplus \bZ_{+}^{r_+} \oplus \bZ_{-}^{r_-} \cong \bZ_d^{s_d} \oplus \bZ_{+}^{s_+} \oplus \bZ_{-}^{s_-}
\]
for $r_d, r_+, r_-, s_d, s_+, s_- \in \bZ_{\geq 0}$. Applying $-\otimes_{\bZ} \bQ$, $r_+ + r_d = s_+ + s_d$ and $r_- + r_d = s_- + s_d$, and applying $-\otimes_{\bZ} \bF_2$, $r_d = s_d$, as $\bZ_+$ and $\bZ_-$ reduce to the trivial $\bF_2[C_2]$-module, and $\bZ_d \cong \bZ[C_2]$ reduces to the indecomposable module $\bF_2[C_2]$. Therefore $r_+ = s_+$, $r_- = s_-$, and $r_d = s_d$.
\end{proof}

\subsection{Levy Structures}

We now introduce the notion of a Levy structure, which will be used below to understand general $C_2$-modules, which aren't necessarily torsion-free.

Finitely generated $\bZ [C_p]$-modules, where $p$ is prime, were classified by Lawrence Levy \cite{Levy81,Levy85,Levy85a}.
For $p=2$, his classification relies on and is stated in terms of the presentation 
\[ \bZ [C_2] \ \cong \ 
\bZ \times_{\bF_2} \bZ \  = \
\{ (x,y) \in \bZ \times \bZ \,\mid\, x+y \mbox{ is even}\}.
\]
Whilst our classification method is different and more explicit, the core of both methods is to reduce the classification
to some linear algebraic data. To acknowledge Levy's contribution, we make the following definition.

\begin{defn}
A Levy structure is a triple $Q=(B, V_+, V_-, {\theta_+},{\theta_-})$, where
\begin{itemize}
  \item $B$ is a finite $C_2$-module,
  \item $V_+, V_-$ are finite-dimensional $\bF_2$-vector spaces,
  \item $\theta_+ \colon V_+ \rightarrow \HH^1(B)$, $\theta_- \colon V_- \rightarrow \HH^2(B)$ are $\bF_2$-linear maps.
\end{itemize}
\end{defn}

Levy structures form a category $\sL$. The morphisms are triples
\[f=(f_0,f_+,f_-): (B, V_+, V_-, {\theta_+}, {\theta_-} ) \to (B', V_+', V_-', {\eta_+}, {\eta_-} ),\]
where $f_0: B\to B'$ is a homomorphism of $C_2$-modules and $f_+ : V_+ \to V_+'$,
$f_- : V_- \to V_-'$ are $\bF_2$-linear maps, making the following diagrams commutative:
\begin{equation} \label{equality_1}
  \begin{tikzcd}
	V_+ & \HH^1(B)  \\
	V_+' & \HH^1(B') 
	\arrow["\theta_+"', from=1-1, to=1-2]
	\arrow["H^1(f_0)", from=1-2, to=2-2]
	\arrow["\eta_+", from=2-1, to=2-2]
	\arrow["f_+"', from=1-1, to=2-1]
  \end{tikzcd}
\quad \quad
  \begin{tikzcd}
	V_- & \HH^2(B)  \\
	V_-' & \HH^2(B') 
	\arrow["\theta_-"', from=1-1, to=1-2]
	\arrow["H^2(f_0)", from=1-2, to=2-2]
	\arrow["\eta_-", from=2-1, to=2-2]
	\arrow["f_-"', from=1-1, to=2-1]
  \end{tikzcd}
\end{equation}
The category 
$\sL$ is additive and has direct sums
\[(B, V_+, V_-, {\theta_+}, {\theta_-} ) \oplus (B', V_+', V_-', {\eta_+}, {\eta_-} ) \coloneqq (B\oplus B', V_+ \oplus V_+', V_- \oplus V_-', {\theta_+}\oplus{\eta_+}, {\theta_-}\oplus{\eta_-} )\]
so that we can talk about indecomposable objects. Moreover, $\sL$ has split idempotents, i.e., for $Q \in \sL$, every $e=e^2 \in \End_{\sL}(Q)$ yields a direct sum decomposition of $Q$. 
\begin{prop} \label{sLk}
  The category $\sL$ is a Krull-Schmidt category.
\end{prop}

  Recall that a Krull-Schmidt category is an additive category, where every object decomposes into a finite direct sum of objects
  with local endomorphism rings \cite{Krause}.

\begin{proof}
Given the finite nature of  $\sL$, every object there decomposes into a direct sum of indecomposable objects. Suppose $X\in\sL$ is indecomposable.
  Then the ring $\End(Q) = \sL (Q,Q)$ has no non-trivial idempotents.
  On the other hand, $\End(Q)$ is a finite ring. It remains to observe that a finite ring $R$ without non-trivial idempotents is necessarily local.

  Recall that a ring $R$ is local if the sum of two non-units is again non-unit.
  The Jacobson radical $J(R)$ is nilpotent since $R$ is finite.
  Hence, it suffices to observe that $R/J(R)$ is a division ring. Suppose not. Then it has a non-trivial idempotent $e\in R/J(R)$, $e\not\in\{0,1\}$.
  But an idempotent can be lifted modulo a nilpotent ideal. This gives a non-trivial idempotent $e_1 \in R$.
\end{proof}  

\begin{cor} \cite[Theorem 4.2]{Krause}\label{cor:KSforLevyStr}
  Consider $X\in \sL$ and two decompositions
  $X= X_1 \oplus \ldots \oplus X_r= Y_1 \oplus \ldots \oplus Y_s$ into non-zero indecomposable objects.
  Then $r=s$ and there exists a permutation $\sigma$ such that $X_i \cong Y_{\sigma (i)}$ for all $i$.
\end{cor}

We can classify indecomposable Levy structures into three types.

\begin{lemma}\label{lem:indecomlevystr}
  Any indecomposable object $X \in \sL$ is of one of the following types:
  \begin{enumerate}
    \item $X \cong (B,0,0,0,0)$, where $B = (\bZ / p^r \bZ)_+$ or $B = (\bZ / p^r \bZ)_-$, for $r \geq 1$ and $p$ an odd prime.
    \item $X \cong (0,\bF_2,0,0,0)$ or $X = (0, 0, \bF_2,0,0)$.
    \item $X \cong (B, V_+, V_-, {\theta_+},{\theta_-})$, indecomposable, where $B$ is $2$-power-torsion and non-zero.
  \end{enumerate}
\end{lemma}

\begin{proof}
  Let $X = (B, V_+, V_-, {\theta_+},{\theta_-})$ be indecomposable, and first suppose that $B_{(p)} \neq 0$ for some odd prime $p$. Then $\HH^1(B_{(p)}) = 0$ by Lemma \ref{cohom_div_by2}, and so $\theta_+ \colon V_+ \rightarrow \HH^1(B)$ is valued in 
  \[
  \HH^1(\oplus_{l \neq p} B_{(l)}) \subset \HH^1(\oplus_{l \neq p} B_{(l)}) \oplus \HH^1(B_{(p)}) = \HH^1(B).
  \]
  Similarly for $\HH^2(B_{(p)})$ and $\theta_- \colon V_- \rightarrow \HH^2(B)$, and therefore 
  \[
    (B_{(p)}, 0, 0, 0, 0) \oplus (\oplus_{l \neq p} B_{(l)}, V_+, V_-, \theta_+, \theta_-) \cong (B, V_+, V_-, {\theta_+},{\theta_-})
  \]
  is a direct sum decomposition in $\sL$. As $B_{(p)} \neq 0$, this forces $B_{(l)} = 0$ for any prime $l \neq p$, and therefore $B_{(p)} = B$. Now $\HH^1(B) = 0$ and $\HH^2(B) = 0$, again by Lemma \ref{cohom_div_by2}, so $\theta_+ = 0$, $\theta_- = 0$, and
  \[
  (B_{(p)}, 0, 0, 0, 0) \oplus (0, V_+, V_-, 0, 0) \cong (B, V_+, V_-, {\theta_+},{\theta_-})
  \]
  is a direct sum decomposition in $\sL$, and hence $V_+ = 0$ and $V_- = 0$, as $B \neq 0$. Because $B$ admits division by $2$, $B = (B^0)_+ \oplus (B^1)_-$ by Lemma \ref{cohom_div_by2}, and therefore 
  \[
  X = ((B^0)_+,0,0,0,0) \oplus ((B^1)_-,0,0,0,0)
  \]
  which forces $B = A_+$ or $B = A_-$, where $A$ is the underlying abelian group of $B$. In either case, any proper non-trivial decomposition of $A$ as an abelian group induces a proper non-trivial decomposition of $B$ as a $C_2$-module, so $A$ must be indecomposable as an abelian group, hence cyclic of odd prime power order, and we see that $X$ is of type (1).
  
  If instead $B_{(p)} = 0$ for all odd primes, then $B = B_{(2)}$. If $B=0$, then $\theta_+ = 0$ and $\theta_- = 0$, and so any decomposition of either $V_+$ or $V_-$ into proper non-trivial subspaces induces a decomposition of $X$. In particular, $X = (0,\bF_2,0,0,0)$ or $X = (0, 0, \bF_2,0,0)$, by the indecomposability of $X$, and $X$ is of type (2). If instead $B \neq 0$, then $B = B_{(2)}$, and $X$ is of type (3). 
\end{proof}

\subsection{Mixed Modules}

We are now interested in the following class of modules. For a $C_2$-module $M$, we write $M_T$ for the $\bZ$-torsion subgroup, which is always a $C_2$-submodule. By Proposition \ref{prop:torsionfreeKS}, we know that the $C_2$-module $M / M_T$ is isomorphic to a direct sum of copies of $\bZ_+$, $\bZ_-$ and $\bZ_d$.

\begin{lemma}\label{lem:mixedequiv}
  Let $M$ be a finitely generated $C_2$-module. The following are equivalent:
  \begin{enumerate}
    \item $M / M_T \cong \bZ_+^n \oplus \bZ^m_-$ for some $m,n \geq 0$,
    \item $M$ is an extension of $\bZ_+^n \oplus \bZ^m_-$ by a finite $C_2$-module, for some $m,n \geq 0$,
    \item $M \cong B_{\Theta}(\bZ)$ for some $\Theta$ and finite $C_2$-module $B$.
  \end{enumerate}
\end{lemma}

\begin{proof}
  The fact that (1) implies (2) follows from the exact sequence $0 \rightarrow M_T \rightarrow M \rightarrow M / M_T \rightarrow 0$, and the fact that (2) implies (3) from Corollary \ref{even_part_ext_new}. The fact that (3) implies (1) is direct from the definition of the $C_2$-module $B_{\Theta}(\bZ)$.
\end{proof}

\begin{defn}
A \emph{mixed module} is a finitely generated $C_2$-module which satisfies the equivalent conditions of Lemma \ref{lem:mixedequiv}
\end{defn}

Let $\sM$ denote the full subcategory of $C_2$-modules, with objects the mixed modules. We now define a functor
\[
\sL(-) \colon \sM \rightarrow \sL
\]
to the category of Levy structures, which will use to deduce properties of $\sM$ from properties of $\sL$.

Recall the invariants $(-)^0$ and skew-invariants $(-)^1$, as defined at the start of Section \ref{sect:cohomC2modules}.
\begin{defn}
  For $M \in \sM$, we set
  \[
  M(0) \coloneqq (M / M_T)^0 \in \Ab, \qquad M(1) \coloneqq (M / M_T)^1 \in \Ab.
  \]
\end{defn}

\begin{remark}
  We note that for any $M \in \sM$, $M / M_T = M(0) \oplus M(1)$ as abelian groups.
\end{remark}

To define $\sL(-)\colon \sM \rightarrow \sL$ on objects, let $M$ be a mixed module, and let $s_M \colon M / M_T \rightarrow M$ be a splitting of $\pi_M \colon M \rightarrow M / M_T$. We note that this is equivalent to the data of an abelian group isomorphism $\gamma_M \colon M \xrightarrow{\sim} M_T \oplus M / M_T$ such that the diagram
% https://q.uiver.app/#q=WzAsMTAsWzIsMCwiTSJdLFsyLDEsIk1fVCBcXG9wbHVzIE0gLyBNX1QiXSxbMywxLCJNIC8gTV9UIl0sWzMsMCwiTSAvIE1fVCJdLFsxLDAsIk1fVCJdLFsxLDEsIk1fVCJdLFs0LDEsIjAiXSxbNCwwLCIwIl0sWzAsMCwiMCJdLFswLDEsIjAiXSxbMCwxLCJcXGdhbW1hIl0sWzQsNSwiXFxpZCJdLFs1LDFdLFs0LDBdLFswLDNdLFsxLDJdLFsyLDZdLFs5LDVdLFs4LDRdLFszLDIsIlxcaWQiXV0=
\[\begin{tikzcd}
	0 & {M_T} & M & {M / M_T} & 0 \\
	0 & {M_T} & {M_T \oplus M / M_T} & {M / M_T} & 0
	\arrow[from=1-1, to=1-2]
	\arrow[from=1-2, to=1-3]
	\arrow["\id", from=1-2, to=2-2]
	\arrow[from=1-3, to=1-4]
	\arrow[from=1-4, to=1-5]
	\arrow["\gamma_M", from=1-3, to=2-3]
	\arrow["\id", from=1-4, to=2-4]
	\arrow[from=2-1, to=2-2]
	\arrow[from=2-2, to=2-3]
	\arrow[from=2-3, to=2-4]
	\arrow[from=2-4, to=2-5]
\end{tikzcd}\]
commutes.

\begin{defn}
  For $M \in \sM$, we set
  \[
  \overline{M(0)} \coloneqq M(0) / 2 M(0) \in \Vect_{\bF_2}, \qquad \overline{M(1)} \coloneqq M(1) / 2 M(1) \in~\Vect_{\bF_2}.
  \]
\end{defn}

Using a fixed splitting $s_M$, we can make the following definition.

\begin{defn}
  For $M \in \sM$, we set
  \[
  \sL(M) \coloneqq (M_T, \overline{M(0)}, \overline{M(1)}, \theta_+, \theta_-),
  \]
  where
  \begin{align*}
    \widetilde{\theta}_+ \colon M(0) \rightarrow M_T^1, \qquad \widetilde{\theta}_+ = \pi_{M_T} \circ \gamma_M \circ \epsilon_M \circ s_M, \\
    \widetilde{\theta}_- \colon M(1) \rightarrow M_T^0, \qquad \widetilde{\theta}_- = \pi_{M_T} \circ \gamma_M \circ \epsilon_M \circ s_M,
  \end{align*}
  and
  \begin{align*}
    \theta_+ \colon \overline{M(0)} \rightarrow M_T^1 / (\epsilon - 1)M_T = \HH^1(M_T), \\
    \theta_- \colon \overline{M(1)} \rightarrow M_T^0 / (\epsilon + 1)M_T = \HH^2(M_T),
  \end{align*}
  are the induced maps. If $f \colon M \rightarrow N$ is a morphism in $\sM$, we set
  \[
  \sL(f) = (f_T, f(0), f(1)),
  \]
  where
  \[
    f_T \colon M_T \rightarrow N_T, \qquad
    f(0) \colon \overline{M(0)} \rightarrow \overline{N(0)}, \qquad
    f(1) \colon \overline{M(1)} \rightarrow \overline{N(1)},
  \]
  are the natural maps induced by $f$.
\end{defn}

Let us verify that this is well-defined.

\begin{lemma}
  For $M \in \sM$, the Levy structure $\sL(M)$ is independent of the choice of splitting $s_M$, and $\sL(f) \colon \sL(M) \rightarrow \sL(N)$ is a well-defined morphism in $\sL$.
\end{lemma}

\begin{proof}
  Suppose that $s_M, s_M' \colon M / M_T \rightarrow M$ are two choices of splitting, with corresponding isomorphisms $\gamma_M, \gamma_M' \colon M \xrightarrow{\sim} M_T \oplus M/ M_T$ respectively. An alternative description of the maps $\widetilde{\theta}_+, \widetilde{\theta}_+' \colon M(0) \rightarrow M_T^1$, are that $\widetilde{\theta}_+(x) = \pi_{M_T} (\epsilon_*(0,x))$ and $\widetilde{\theta}_+'(x) = \pi_{M_T} (\epsilon_*'(0,x))$, where $\epsilon_* = \gamma_M \circ \epsilon_M \circ \gamma_M^{-1}$ and $\epsilon_*' = \gamma_M' \circ \epsilon_M \circ (\gamma_M')^{-1}$ are the transported actions of $\epsilon \in \bZ[C_2]$ on $M_T \oplus M / M_T$. Because these are defined to make $\gamma_M$ and $\gamma_M'$ both $\bZ[C_2$]-equivariant, we have that $\phi \circ \epsilon_* = \epsilon_*' \circ \phi$, where $\phi \coloneqq \gamma_M' \circ \gamma_M^{-1}$. Now $\phi$ is an abelian group automorphism of $M_T \oplus M / M_T$ which is the identity on $M_T$ and the identity after projecting to $M / M_T$, are therefore $\phi$ and $\phi^{-1}$ have the form
  \[
  \phi(m, x) = (m + u(x), x), \qquad \phi^{-1}(m,x) = (m - u(x), x)
  \]
  for some $u \in \Hom_{\bZ}(M/M_T, M_T)$. We can therefore compute that for $x \in M(0)$, 
  \begin{align*}
    \widetilde{\theta}_+(x) &= \pi_{M_T} (\epsilon_*(0,x)), \\
    &= \pi_{M_T}((\phi^{-1} \circ \epsilon_*' \circ \phi)(0,x)), \\
    &= \pi_{M_T}(\phi^{-1} (\epsilon_*'(u(x), x))), \\
    &= \pi_{M_T}(\phi^{-1} (\epsilon_M(u(x)) + \widetilde{\theta}_+'(x), x)), \\
    &= \pi_{M_T}(\epsilon_M(u(x)) + \widetilde{\theta}_+'(x) - u(x), x), \\
    &= \widetilde{\theta}_+'(x) + (\epsilon_M - 1)u(x),
  \end{align*}
  and therefore the induced maps $\theta_+$ and $\theta_+'$ are equal. Taking instead $x \in M(1)$, from an analogous computation one obtains that $\theta_- = \theta_-'$. 
  
  For the second claim, given $f \colon M \rightarrow N$ a morphism in $\sM$, to check that $\sL(f) \colon \sL(M) \rightarrow \sL(N)$ is a morphism in $\sL$ one needs to show that $\theta_{+,N} \circ f(0) = \HH^1(f_T) \circ \theta_{+,M}$ and $\theta_{+,N} \circ f(1) = \HH^2(f_T) \circ \theta_{+,M}$. We show the first, the second following mutatis mutandis. The induced map
  \[
  \widetilde{f} \coloneqq \gamma_N \circ f \circ \gamma_M^{-1} \colon M_T \oplus M/M_T \rightarrow N_T \oplus N / N_T
  \]
  is $\bZ[C_2]$-linear when $M_T \oplus M/M_T$ and $N_T \oplus N / N_T$ are given the transported actions $\epsilon_{*,M} = \gamma_M \circ \epsilon_M \circ \gamma_M^{-1}$ and $\epsilon_{*,N} = \gamma_N \circ \epsilon_N \circ \gamma_N^{-1}$ respectively as above. Similarly to above, because $M_T$ is torsion and $N / N_T$ is torsion free, the map $\widetilde{f}$ has the form
  \[
  \widetilde{f}(m,x) = (f_T(m) + u(x), f_{M/M_T}(x))
  \]
  for some $u \in \Hom_{\bZ}(M/M_T, N_T)$. We can then compute that for $x \in M(0)$,
  \[
    (\widetilde{f} \circ \epsilon_{*,M})(0,x) = \widetilde{f}(\widetilde{\theta}_{+, M}(x),x) = (f_T(\widetilde{\theta}_{+, M}(x)) + u(x), f_{M/M_T}(x)),
  \]
  whereas on the other hand
  \[
  (\epsilon_{*,N} \circ \widetilde{f})(0,x) = \epsilon_{*,N}(u(x), f_{M/M_T}(x)) = (\epsilon_N(u(x)) + \widetilde{\theta}_{+,N}(f_{M/M_T}(x)), f_{M/M_T}(x)).
  \]
  Then comparing the first coefficient,
  \[
  f_T(\widetilde{\theta}_{+, M}(x)) - \widetilde{\theta}_{+,N}(f_{M/M_T}(x))  = (\epsilon_N - 1)(u(x))
  \]
  which when taking the induced map on the quotient implies that $\HH^1(f_T) \circ \theta_{+,M} - \theta_{+,N} \circ f(0) = 0$.
\end{proof}

\begin{eg}\label{eg:stdexample}
Suppose that $M = B_{\Theta}(\bZ)$, where $M / M_T$ is canonically identified with $\bZ_+^n \oplus \bZ_-^m$, and take $s_M \colon \bZ_+^n \oplus \bZ_-^m \rightarrow B_{\Theta}(\bZ)$ to be the standard splitting. Then $M(0) = \bZ_+^n$, $M(1) = \bZ_-^m$, which induce identifications $\overline{M(0)} = \bF_2^n$ and $\overline{M(1)} = \bF_2^m$. The maps $\theta_+ \colon \bF_2^n \rightarrow \HH^1(B)$ and $\theta_- \colon \bF_2^m \rightarrow \HH^2(B)$ are simply those induced by $\Theta_+ \colon \bZ^n \rightarrow B^1$ and $\Theta_- \colon \bZ^m \rightarrow B^0$.
\end{eg}

We can now describe how $\sM$ and $\sL$ are related.

\begin{prop}\label{prop:Lfunctorproperties}
  The functor $\sL(-) \colon \sM \rightarrow \sL$ satisfies the following properties.
  \begin{enumerate}
    \item $\sL(-) \colon \sM \rightarrow \sL$ is additive, commutes with direct sums, and is essentially surjective.
    \item If $M, N \in \sM$, then $M \cong N$ in $\sM$ if and only if $\sL(M) \cong \sL(N)$ in $\sL$.
    \item $M$ is indecomposable in $\sM$ if and only if $\sL(M)$ is indecomposable in $\sL$.
  \end{enumerate}
\end{prop}

\begin{remark}\label{rem:notionsofdecomp}
  Note that $\sM$ is closed under direct summands in the category of all $C_2$-modules, and therefore indecomposability in $\sM$ is equivalent to indecomposability as a $C_2$-module. Indeed, if $M \in \sM$ and $M = M_1 \oplus M_2$ is a decomposition as $C_2$-modules, then $M_1 / M_{1, T} \oplus M_2 / M_{2, T} \cong M / M_T$, and therefore both $M_1$ and $M_2$ are mixed modules, by Proposition \ref{prop:torsionfreeKS}.
\end{remark}

\begin{proof}
  It is clear that $\sL(-)$ is additive and commutes with direct sums. For the essential surjectivity, let $(B, V_+, V_-, \theta_+, \theta_-) \in \sL$, and choose a basis for $V_+$ and $V_-$. Any lift of the maps $\theta_+ \colon \bF_2^n \rightarrow \HH^1(B)$ and $\theta_- \colon \bF_2^m \rightarrow \HH^2(B)$ to $\Theta_+ \colon \bZ^n \rightarrow B^1$ and $\Theta_- \colon \bZ^m \rightarrow B^0$ respectively allow us to define $B_{\Theta}(\bZ) \in \sM$, with the property that $\sL(\sM) = (B, V_+, V_-, \theta_+, \theta_-)$, by Example \ref{eg:stdexample}.

  For the second point, note that by functoriality $\sL(M) \cong \sL(N)$ whenever $M \cong N$, so suppose conversely that $\sL(M) \cong \sL(N)$. Then $M_T \cong N_T$, $\dim_{\bF_2}\overline{M(0)} = \dim_{\bF_2}\overline{M(0)}$ and $\dim_{\bF_2}\overline{M(0)} = \dim_{\bF_2}\overline{M(0)}$, and we may therefore pick isomorphisms $M \cong B_{\Theta}(\bZ)$, $N \cong B_{\Psi}(\bZ)$ for the same finite $C_2$-module $B$ and $\Theta_+, \Psi_+ \colon \bZ^n \rightarrow B$, $\Theta_-, \Psi_- \colon \bZ^m \rightarrow B$ with the same domains. Fixing an isomorphism $(f_0,f_+,f_-) \colon \sL(B_{\Theta}(\bZ)) \rightarrow \sL(B_{\Psi}(\bZ))$, we may lift the isomorphisms $f_+ \colon \bF_2^n \rightarrow \bF_2^n$ and $f_- \colon \bF_2^m \rightarrow \bF_2^m$ to isomorphisms $S_+ \colon \bZ^n \rightarrow \bZ^n$ and $S_- \colon \bZ^m \rightarrow \bZ^m$, using that $\GL_n(\bZ) \rightarrow \GL_2(\bF_2)$ is surjective, the latter being generated by transvections. 

  Writing $\sL(B_{\Theta}(\bZ)) = (B, \bF_2^n, \bF_2^m, \theta_+, \theta_-)$, $\sL(B_{\Psi}(\bZ)) = (B, \bF_2^n, \bF_2^m, \psi_+, \psi_-)$, we have that 
  \begin{align*}
    \delta_+ &\coloneqq \Psi_+ \circ S_+ - f_0 \circ \Theta_+ \colon \bZ^n \rightarrow B^1, \\
    \delta_- &\coloneqq \Psi_- \circ S_- - f_0 \circ \Theta_- \colon \bZ^m \rightarrow B^0
  \end{align*}
  reduce to
  
  \begin{align*}
    \psi_+ \circ f_+ - \HH^1(f_0) \circ \theta_+ &\colon \bF_2^n \rightarrow \HH^1(B), \\
    \psi_- \circ f_- - \HH^2(f_0) \circ \theta_+ &\colon \bF_2^n \rightarrow \HH^2(B)
  \end{align*}
  respectively, and therefore $\delta_+$ and $\delta_-$ are valued in $(\epsilon - 1)B$ and $(\epsilon + 1)B$ respectively. Defining
  \begin{align*}
    \Psi_+' &\coloneqq \Psi_+ - \delta_+ \circ S_+^{-1} \colon \bZ^n \rightarrow B^1, \\
    \Psi_-' &\coloneqq \Psi_- - \delta_- \circ S_-^{-1} \colon \bZ^m \rightarrow B^0,
  \end{align*}
  we have that $\Psi_+ - \Psi_+'$ is valued in $(\epsilon - 1)B$ and $\Psi_- - \Psi_-'$ is valued in $(\epsilon + 1)B$, hence $B_{\Psi}(\bZ) \cong B_{\Psi'}(\bZ)$, even as extensions, by Corollary \ref{even_part_ext_new}. It is thus sufficient to give an isomorphism $\phi \colon B_{\Theta}(\bZ) \xrightarrow{\sim} B_{\Psi'}(\bZ)$, which we do by taking $\phi$ to be explicitly defined by
  \[
  \phi(b,x,y) \coloneqq (f_0(b),S_+(x), S_-(y)).
  \]
  which is a $C_2$-module isomorphism, using the fact that $\Psi_+' \circ S_+ = f_0 \circ \Theta_+$ and $\Psi_-' \circ S_- = f_0 \circ \Theta_-$.

  For the third and final claim, note that any proper direct sum decomposition of $M$ induces a proper direct sum decomposition of $\sL(M)$. Conversely, if $\sL(M) \cong \sL_1 \oplus \sL_2$ is a proper direct sum decomposition of $\sL(M)$, then taking $M_1, M_2 \in \sM$ with $\sL(M_1) \cong \sL_1$, $\sL(M_2) \cong \sL_2$, which is possible by the essential surjectivity of $\sL(-)$, we have that $\sL(M) \cong \sL(M_1) \oplus \sL(M_2) \cong \sL(M_1 \oplus M_2)$, and therefore, using part (2), $M \cong M_1 \oplus M_2$ is a proper direct sum decomposition of $M$.
\end{proof}

\begin{cor}\label{cor:KSformixedmodules}
Any non-zero mixed module $M$ can be written as a direct sum of indecomposable $C_2$-modules. This direct sum description is unique up to isomorphism and permutation, and each direct-summand is a mixed module.
\end{cor}

\begin{proof}
  For $M \in \sM$, we may write $\sL(M) = \sL_1 \oplus \cdots \oplus \sL_r$ as a direct sum of indecomposable objects of $\sL$ by Corollary \ref{cor:KSforLevyStr}, and lift them to $M_1, ... , M_r \in \sM$ by Proposition \ref{prop:Lfunctorproperties}(1). Because $\sL(-)$ commutes with direct sums, $\sL(M) \cong \sL(M_1) \oplus \cdots \oplus \sL(M_r)$, and therefore $M \cong M_1 \oplus \cdots \oplus M_r$ by Proposition \ref{prop:Lfunctorproperties}(2), with each $M_i$ indecomposable by Proposition \ref{prop:Lfunctorproperties}(3) and Remark \ref{rem:notionsofdecomp}. 

  If $M_1 \oplus \cdots \oplus M_r \cong M \cong N_1 \oplus \cdots \oplus N_s$ are two indecomposable decompositions of $M$, then $\sL(M_1) \oplus \cdots \oplus \sL(M_r) \cong \sL(M) \cong \sL(N_1) \oplus \cdots \oplus \sL(N_r)$ are indecomposable decompositions in $\sL$, by Proposition \ref{prop:Lfunctorproperties}(3), and therefore the uniqueness follows from the uniqueness of Corollary \ref{cor:KSforLevyStr} and Proposition \ref{prop:Lfunctorproperties}(2). The final claim follows from Remark \ref{rem:notionsofdecomp}.
\end{proof}


\subsection{General classification} We can now state our classification of all $C_2$-modules.

\begin{thm}\label{thm:generalclassification} 
  Any finitely generated $C_2$-module is a direct sum of indecomposable finitely generated $C_2$-modules, and this direct sum decomposition is unique up to isomorphism and permutation of the factors. The indecomposable $C_2$-modules fall into one of the following types.
\begin{enumerate}
\item $p$-Torsion: $B_+$ and $B_-$, where $B$ is a cyclic group of order $p^n$, $n \geq 1$, $p$ odd.
%\item $B_-$, where $B$ is a cyclic group of order $p^n$, $p$ is odd.
\item Torsion-free: $\bZ_d$, $\bZ_+$ and $\bZ_-$.
\item Mixed: $B_\Theta(\bZ)$ for a $2$-power-torsion $C_2$-module $B \neq 0$, with $\sL(B_\Theta(\bZ))$ indecomposable.
\end{enumerate}
\end{thm}

\begin{remark}
  Note that the isomorphism class of the modules $B_\Theta(\bZ)$ depends only on the isomorphism class of the associated Levy structure $\sL(B_\Theta(\bZ))$ in $\sL$, by Proposition \ref{prop:Lfunctorproperties}(2).
\end{remark}

\begin{proof}
  Suppose that $A$ is a finitely generated $C_2$-module. From Proposition \ref{prop:torsionfreeKS}, $A/A_T \cong \bZ_+^n \oplus \bZ_-^m \oplus \bZ_d^r$ for some $r,m,n \geq 0$. The module $\bZ_d \cong \bZ[C_2]$ is free, and therefore we may split the projection $A \twoheadrightarrow \bZ_d^r$ to write $A \cong M \oplus \bZ_d^r$ for some $C_2$-submodule $M$ of $A$. We have that $M / M_T \cong \bZ_+^n \oplus \bZ_-^m$, and therefore $M$ is a mixed module. In particular, $M$ decomposes into indecomposable $C_2$-modules by Corollary \ref{cor:KSformixedmodules}, and therefore so does $A$. 

  To see that all indecomposable $C_2$-modules are included in our list, and all modules on our list our indecomposable, note that our list consists of exactly the $C_2$-module $\bZ_d$, indecomposable by Proposition \ref{prop:torsionfreeKS}, and the indecomposable $C_2$-modules contained in $\sM$, by Lemma \ref{lem:mixedequiv}, Proposition \ref{prop:Lfunctorproperties}(3) and Remark \ref{rem:notionsofdecomp}, as these are exactly the mixed modules reducing to the indecomposable Levy structures of Lemma \ref{lem:indecomlevystr}. To see that any finitely generated indecomposable $C_2$-module is either $\bZ_d$ or contained in $\sM$, note that if $A$ is a finitely generated indecomposable $C_2$-module, then we can write $A \cong M \oplus \bZ_d^r$ as above. If $M = 0$, then $r = 1$ and $A = \bZ_d$ by Proposition \ref{prop:torsionfreeKS}, and if $M \neq 0$, then $r = 0$ and $A = M$ is a mixed module.

  Finally, to show that the decomposition into indecomposables is unique, suppose that we have two descriptions of $A$ as a direct sum of indecomposable modules on our list. Grouping together the terms other than $\bZ_d$, we therefore have an isomorphism
  \[
    M \oplus \bZ_d^r \cong A \cong N \oplus \bZ_d^s
  \]
for mixed modules $M,N$ and $r,s \geq 0$. From the uniqueness of the decomposition of the torsion-free $C_2$-module $A / A_T$ (Proposition \ref{prop:torsionfreeKS}), we have that $r = s$, and therefore
\[
M^{\oplus n} \cong N^{\oplus n}
\]
for some $n \geq 1$ because $\bZ[C_2]$, the endomorphism ring of $\bZ_d$, is a commutative $\bZ$-algebra which is finite-free over $\bZ$ \cite[Cor.\ 2.2, Thm.\ 3.10]{GOO}. Then writing $M$ and $N$ as their indecomposable decompositions, these are the same up to isomorphism and permutation by Corollary \ref{cor:KSformixedmodules} as $M$ and $N$ are mixed, and therefore the same is true of the two indecomposable decompositions of $A$.
\end{proof}

In light of Theorem \ref{thm:generalclassification}, to understand general $C_2$-modules one is therefore reduced to understanding when the $C_2$-module $B_{\Theta}(\bZ)$ is indecomposable for a finite non-trivial $2$-power-torsion $C_2$-module $B$. From Proposition \ref{prop:Lfunctorproperties}(3) this is equivalent to the following problem.

\begin{prob}
Classify indecomposable Levy structures $(B, V_+, V_-, \theta_+, \theta_-)$ for which $B$ is a finite non-trivial $2$-power-torsion $C_2$-module.
\end{prob}

For example, is it possible to characterise the indecomposability of $(B, V_+, V_-, \theta_+, \theta_-)$ in terms of the indecomposability of $B$ and conditions on the maps $\theta_+$ and $\theta_-$?
